\section{动机}

相机标定是三维计算机视觉中从二维图像提取度量信息的必要步骤。该领域已经开展了大量工作，起初来自摄影测量学界（例如见文献~\cite{brown1971close,faig1975calibration}），近来则更多来自计算机视觉领域（例如见文献~\cite{gennery1979stereo,ganapathy1984decomposition,tsai1987versatile,faugeras1986stereo,weng1992camera,wei1993complete,maybank1992theory,faugeras1992self,triggs1998autocalibration}）。粗略地说，我们可以将这些技术分为两类：

\textbf{摄影测量标定。} 标定通过观测一个标定物体来完成，该物体在三维空间中的几何形状已知且具有很高的精度。标定可以非常高效地完成~\cite{faugeras1993three}。标定物体通常由两个或三个彼此正交的平面组成。有时也使用一个经历精确已知平移的平面~\cite{tsai1987versatile}。这些方法需要昂贵的标定装置和复杂的设置。

\textbf{自标定。} 这类技术不使用任何标定物体。仅通过在静态场景中移动相机，场景的刚性通常就可以利用纯图像信息，从一次相机位移中为相机的内参数提供两个约束~\cite{maybank1992theory}。因此，如果图像由同一台具有固定内参数的相机拍摄，那么三幅图像之间的对应关系就足以恢复内参数和外参数，并据此在相似变换意义下重建三维结构~\cite{luong1997self,hartley1994several}。尽管这种方法非常灵活，但目前还不够成熟~\cite{bougnoux1998projective}。由于需要估计的参数很多，我们无法始终获得可靠的结果。

还存在其他技术，例如利用正交方向的消失点进行标定~\cite{caprile1990vanishing,liebowitz1998metric}，以及利用纯旋转进行标定~\cite{hartley1994rotating,stein1995accurate}。

我们当前的研究聚焦于桌面视觉系统（desktop vision system，DVS），因为使用 DVS 的潜力很大。相机正变得廉价且无处不在。DVS 面向普通大众，而他们并非计算机视觉专家。典型的计算机用户只会偶尔执行视觉任务，因此不愿意为昂贵的设备投资。因此，灵活性、鲁棒性和低成本十分重要。本文所述的相机标定技术正是在这些考虑下开发的。

所提出的技术仅要求相机在少数（至少两个）不同方向下观测一个平面模板。该模板可以用激光打印机打印出来，并附着在一个“合理的”平面表面上（例如硬质书封面）。相机或平面模板均可用手移动，且不需要知道其运动。所提出的方法介于摄影测量标定和自标定之间，这是因为我们使用的是二维度量信息，而不是三维度量信息或纯粹的隐式信息。我们使用计算机仿真数据和真实数据对所提出的技术进行了测试，并获得了非常好的结果。与传统方法相比，所提出的技术灵活得多；与自标定相比，它具有相当程度的鲁棒性。我们相信，这项新技术使三维计算机视觉从实验室环境向现实世界迈进一步。

值得注意的是，Bill Triggs~\cite{triggs1998autocalibration} 最近提出了一种自标定技术，它至少需要对平面场景进行五次观测。该技术比我们的方法更加灵活，但存在初始化困难。Liebowitz 和 Zisserman~\cite{liebowitz1998metric} 描述了一种针对平面透视图像的度量校正技术，所使用的度量信息包括一个已知角度、两个相等但未知的角度以及一个已知长度比。他们还提到，只要存在这样的度量校正平面，就可以对相机的内参数进行标定，尽管他们没有给出算法或实验结果。

本文的组织结构如下。第 2 节描述从观测单个平面得到的基本约束。第 3 节描述标定过程：我们先给出一个闭式解，然后进行非线性优化，同时对镜头径向畸变进行建模。第 4 节研究所提出的相机标定技术失效的配置；在实践中，这些情况很容易避免。第 5 节给出实验结果，并使用计算机仿真数据和真实数据验证所提出的技术。附录给出若干细节，包括估计模型平面与其图像之间单应性的技术。
